\begin{IEEEbiography}[{\includegraphics[width=1in,height=1.25in,clip, keepaspectratio]{sections/a1.jpg}}]{Edara Sreenivasa Reddy} is currently a Professor (HAG) with the School of Computer Science and Engineering, VIT-AP University, Amaravati, India. He has more than 34 years of academic experience and 18 years of research experience. During his academic
career, he has served in several teaching and administrative positions, including Assistant Professor, Associate Professor, Professor, Senior Professor, Principal, Dean, and Vice-Chancellor.

He has published more than 125 Scopus-indexed journal articles, authored 12 books, contributed 12 chapters to edited research volumes, and obtained 12 patents. Under his supervision, 46 Ph.D. scholars and two M.Phil. scholars have completed their degrees. His research interests include artificial intelligence, machine learning, deep learning, generative artificial intelligence, large language models, and quantum computing.

Prof. Reddy is an NVIDIA University Ambassador for
deep learning and an NVIDIA-certified coach for generative artificial intelligence and large language models. He has also completed professional certification in quantum computing from IBM. He received the Best Computer Science Teacher Award from the Indian Society for
Technical Education. He is a Fellow of the Computer Science Research Council, USA, and an Honorary Fellow of the Academy of Science, Technology and Communications. He is a Senior Member of IEEE and a member of the Indian Society for Technical Education.

\end{IEEEbiography}

\begin{IEEEbiography}[{\includegraphics[width=1in,height=1.25in,clip, keepaspectratio]{sections/a2.jpg}}]{Shanmukesh Bonala} is currently pursuing the B.Tech. degree in computer
science and engineering with VIT-AP University,
Amaravati, India. His research interests include artificial intelligence, computer vision, deep learning, model optimization, and energy-efficient artificial intelligence.

He has contributed to research and development projects in artificial intelligence and computer vision and has published research datasets through IEEE DataPort. His work focuses on the development of scalable and reproducible artificial-intelligence systems. He is an IEEE
Student Member.
\end{IEEEbiography}
